\documentclass[tikz,border=8pt]{standalone}
\usepackage{ctex}
\usepackage{amssymb}
\usepackage{pgfplots}
\pgfplotsset{compat=1.18}
\usetikzlibrary{calc,positioning,arrows.meta}

\begin{document}
\begin{tikzpicture}

%% ===== 左图：函数曲线 + 一阶切线 =====
\begin{scope}[xshift=0cm]

  % 坐标轴
  \draw[->, gray!60, thin] (-0.5, 0) -- (5.5, 0) node[right, font=\scriptsize] {$x$};
  \draw[->, gray!60, thin] (0, -0.5) -- (0, 5.0) node[above, font=\scriptsize] {$f(x)$};

  % 函数曲线 f(x) = 0.2*(x-2)^2 + 0.5
  \draw[blue!70!black, thick, domain=0.2:4.8, samples=60]
    plot (\x, {0.2*(\x-2)*(\x-2) + 0.5});
  \node[blue!70!black, font=\small] at (4.5, 3.5) {$f(x)$};

  % 展开点 x0 = 1.5, f(x0) = 0.2*0.25+0.5 = 0.55
  \coordinate (X0) at (1.5, 0.55);
  \fill (X0) circle (2pt);
  \node[below, font=\scriptsize] at (1.5, -0.1) {$x_0$};

  % 一阶切线（斜率 = f'(x0) = 0.4*(x0-2) = -0.2）
  % 切线：y - 0.55 = -0.2*(x - 1.5)  =>  y = -0.2x + 0.3 + 0.55 = -0.2x + 0.85
  \draw[red!80!black, thick, domain=0.3:3.5]
    plot (\x, {-0.2*\x + 0.85});
  \node[red!80!black, font=\scriptsize] at (3.3, 0.25) {一阶切线};

  % 二阶近似抛物线（以 x0=1.5 为中心，系数 0.2）
  % q(x) = f(x0) + f'(x0)*(x-x0) + 0.5*f''(x0)*(x-x0)^2
  %      = 0.55 - 0.2*(x-1.5) + 0.5*0.4*(x-1.5)^2
  %      = 0.55 - 0.2*(x-1.5) + 0.2*(x-1.5)^2
  \draw[green!60!black, thick, dashed, domain=0.2:3.5, samples=40]
    plot (\x, {0.55 - 0.2*(\x-1.5) + 0.2*(\x-1.5)*(\x-1.5)});
  \node[green!60!black, font=\scriptsize] at (0.6, 2.2) {二阶近似};

  % 标注展开点竖虚线
  \draw[gray, dashed, thin] (1.5, 0) -- (X0);

  % 注释框（放在 axis 外）
  \node[draw=gray!50, fill=white, thin, rounded corners=2pt,
        font=\scriptsize, inner sep=3pt] at (4.0, 1.5)
    {$f(x_0+h)\approx f(x_0)$\\$+f'(x_0)h+\frac{f''(x_0)}{2}h^2$};

  % 标题
  \node[font=\small\bfseries] at (2.5, -0.9) {一维：切线（1阶）与二次近似（2阶）};

\end{scope}

%% ===== 右图：多维 Hessian 几何 =====
\begin{scope}[xshift=8.5cm]

  % 模拟二元二次函数 f(x1,x2) = x1^2 + 3*x2^2 的等高线（椭圆）
  % 坐标轴
  \draw[->, gray!60, thin] (-3.5, 0) -- (3.5, 0) node[right, font=\scriptsize] {$x_1$};
  \draw[->, gray!60, thin] (0, -2.8) -- (0, 2.8) node[above, font=\scriptsize] {$x_2$};

  % 等高线（椭圆）
  \draw[gray!40, thin] (0,0) ellipse (0.60 and 0.346);
  \draw[gray!40, thin] (0,0) ellipse (1.20 and 0.693);
  \draw[gray!40, thin] (0,0) ellipse (1.80 and 1.039);
  \draw[gray!40, thin] (0,0) ellipse (2.40 and 1.386);

  % 展开点
  \coordinate (P) at (1.5, 0.7);
  \fill (P) circle (2pt);
  \node[above right, font=\scriptsize] at (P) {$\mathbf{x}_0$};

  % 梯度方向
  % 梯度 = (2*1.5, 6*0.7) = (3, 4.2)，单位化后 * 0.7
  \draw[->, thick, red!80!black, >=Stealth]
    (P) -- ++(0.49, 0.69) node[right, font=\scriptsize] {$\nabla f$};

  % 展开点切平面（用一个矩形模拟切面）
  \draw[orange!70!black, thick, dashed]
    ($(P)+(-1.0,-0.2)$) -- ($(P)+(1.0,0.4)$);
  \node[orange!70!black, font=\scriptsize] at ($(P)+(-0.8,0.4)$) {切平面};

  % Hessian 框（二阶近似信息）
  \node[draw=gray!50, fill=white, thin, rounded corners=2pt,
        font=\scriptsize, inner sep=4pt] at (-1.5, 1.8)
    {$H_f = \nabla^2 f$，对称矩阵，描述曲率};

  % 曲率标注
  \draw[<->, thick, gray!60] (-2.4, -1.385) arc[start angle=210, end angle=330, x radius=2.4, y radius=1.385];
  \node[gray, font=\scriptsize] at (0, -2.1) {曲率 $\kappa\propto\mathbf{v}^\top H\mathbf{v}$};

  % 极小值点标注
  \fill[red!60!black] (0,0) circle (2pt) node[below right, font=\scriptsize] {$\mathbf{x}^*$（极小）};

  % 公式框
  \node[draw=gray!50, fill=white, thin, rounded corners=2pt,
        font=\scriptsize, inner sep=3pt] at (1.0, -2.0)
    {$f(\mathbf{x}_0+\Delta\mathbf{x})\approx f(\mathbf{x}_0)+\nabla f^\top\Delta\mathbf{x}+\frac{1}{2}\Delta\mathbf{x}^\top H\Delta\mathbf{x}$};

  % 标题
  \node[font=\small\bfseries] at (0, -2.8) {多维：Hessian 描述各方向曲率};

\end{scope}

%% 整体标题
\node[font=\large\bfseries] at (8.5, 3.3) {Hessian 二阶 Taylor 展开几何};

\end{tikzpicture}
\end{document}
